% ---------------------------------------------------------------------
% Introduction (~1.5 pages). Numbers marked [DUMMY] are placeholders.
% ---------------------------------------------------------------------

\IEEEPARstart{T}{he} quality of AI-generated speech has crossed a perceptual
threshold. Neural codec language models such as VALL-E~\cite{wang2023valle},
diffusion- and flow-matching-based synthesizers such as StyleTTS~2~\cite{li2023styletts2}
and F5-TTS~\cite{chen2024f5tts}, and massively multilingual zero-shot systems such as
XTTS~\cite{casanova2024xtts}, Seed-TTS~\cite{anastassiou2024seedtts}, and
CosyVoice~\cite{du2024cosyvoice} can now clone a speaker's voice from a few seconds of
reference audio and synthesize speech that human listeners and conventional classifiers
cannot reliably distinguish from genuine recordings. While these systems enable valuable
applications in accessibility, dubbing, and creative production, they have also driven a
sharp rise in documented harms: voice-cloning fraud against individuals and financial
institutions, fabricated audio of political figures circulated during elections, and
evasion of voice-biometric authentication. An analysis of deepfakes actually circulated
online during 2024 found that audio deepfakes are among the fastest-growing and least
reliably detected modalities~\cite{chandra2025deepfakeeval}, and large-scale fake-speech
corpora are now being generated with large language models in the loop to simulate
disinformation campaigns~\cite{luong2024llamapartialspoof}.

The anti-spoofing community has responded with increasingly capable countermeasures.
On the standard ASVspoof~2019 Logical Access (LA) benchmark~\cite{wang2020asvspoof2019},
systems built on self-supervised learning (SSL) speech foundation models with
graph-attention back-ends~\cite{tak2022automatic,jung2022aasist} achieve equal error
rates (EER) below 1\%, and recent state-space-model
detectors~\cite{xiao2025xlsrmamba,chen2024rawbmamba} and lightweight foundation-model
back-ends~\cite{liu2025nes2net} have pushed in-domain accuracy further still. Yet a
growing body of evidence---from the ASVspoof~5 challenge~\cite{wang2024asvspoof5,
wang2025asvspoof5journal} to systematic cross-domain
studies~\cite{muller2022inthewild,muller2024harder}---shows that these in-domain numbers
are a poor predictor of real-world reliability. We identify three failure modes that
jointly block deployment.

\textbf{G1---Generalization to unseen generators.} A detector trained on a fixed set of
text-to-speech (TTS) and voice-conversion (VC) algorithms sees its EER degrade from
below 1\% to 5--30\% when evaluated on synthesizers absent from its training
corpus~\cite{muller2022inthewild,muller2024harder}. Because new synthesis families
appear monthly---neural codec language models and flow-matching synthesizers did not
exist when the dominant training corpora were collected---a practical detector must
generalize \emph{zero-shot}, without retraining, to generator families it has never
seen. The Codecfake study~\cite{xie2025codecfake} demonstrates that detectors trained on
conventional vocoder artifacts largely fail on speech produced through neural audio
codecs, precisely the mechanism used by the newest generation of synthesizers.

\textbf{G2---Channel robustness.} Real-world audio reaches a detector after lossy
compression (MP3, AAC, Opus), telephony coding (AMR, G.711), reverberation, and additive
noise---conditions imposed by messaging platforms, social media transcoding, and phone
networks. The ASVspoof~2021 evaluation~\cite{liu2023asvspoof2021} showed that codec and
compression variability alone can multiply EER several-fold, and the interaction of
channel distortion with spoofing artifacts remains poorly understood. Most published
evaluations still use clean, studio-quality audio.

\textbf{G3---Calibration and explainability.} Detection scores feed downstream
decisions---content moderation, forensic analysis, legal evidence---that require
\emph{calibrated} confidence estimates, not merely a ranking. Recent work shows that
modern SSL-based detectors are systematically over-confident, particularly under domain
shift~\cite{pascu2024calibrated}. Forensic practice additionally demands interpretable
evidence for individual decisions, which binary scores do not provide.

Recent large audio-language models (ALLMs) such as Qwen2-Audio~\cite{chu2024qwen2audio}
and SALMONN~\cite{tang2024salmonn} raise a natural question: does the general-purpose
audio understanding of foundation models subsume the deepfake-detection task? Early
evidence suggests not---prompted zero-shot, ALLMs perform near chance at distinguishing
bona-fide from synthetic speech~\cite{gong2025allm4add}---but no prior work has placed
ALLMs and specialized countermeasures side by side within a controlled cross-domain,
channel-robustness, and calibration protocol. We do so here.

This paper reframes synthetic-speech detection from closed-set binary classification to
\emph{open-set, channel-robust, calibrated} detection, and proposes
\textbf{AuralGuard}, a multi-view one-class detection framework addressing G1--G3
jointly. AuralGuard rests on four design decisions. \emph{First}, it fuses two
complementary views of the input: View~A, a self-supervised foundation-model
representation (WavLM-Large~\cite{chen2022wavlm}) with \emph{learnable layer attention}
that adaptively selects the transformer depths most informative for artifact detection;
and View~B, a hand-crafted feature stream that explicitly exposes low-level vocoder
artifacts through linear-frequency cepstral coefficients (LFCC), modified group delay,
and constant-Q phase---cues that SSL models, pre-trained for semantic tasks, are prone
to discard. \emph{Second}, a gated cross-attention module fuses the two views before a
spectro-temporal graph-attention back-end~\cite{jung2022aasist} produces an
utterance-level embedding. \emph{Third}, a one-class contrastive objective combining
OC-Softmax~\cite{zhang2021oneclass} with a supervised-contrastive
auxiliary~\cite{khosla2020supcon} learns a compact decision region around bona-fide
speech, so that unseen spoofing algorithms are rejected as out-of-distribution rather
than mapped onto decision boundaries fit to known attacks. \emph{Fourth}, a
channel-robustness augmentation curriculum---RawBoost~\cite{tak2022rawboost}, a
realistic codec chain (MP3, AAC, Opus, AMR, G.711), measured room impulse responses,
and additive noise---aligns the training distribution with deployment channels.

We evaluate AuralGuard under a strict zero-shot protocol: all models are trained
\emph{only} on ASVspoof~2019~LA and evaluated without adaptation on six additional
corpora---ASVspoof~2021 LA and DF~\cite{liu2023asvspoof2021},
ASVspoof~5~\cite{wang2025asvspoof5journal}, In-the-Wild~\cite{muller2022inthewild},
WaveFake~\cite{frank2021wavefake}, Codecfake~\cite{xie2025codecfake}, and the
38-language MLAAD corpus~\cite{muller2024mlaad}. Baselines span four generations of
countermeasure design: hand-crafted features with light CNNs, end-to-end raw-waveform
networks, SSL-based detectors including the current published state of the
art~\cite{tak2022automatic,xiao2025xlsrmamba,zhu2024slim}, and zero-shot ALLM
prompting~\cite{gong2025allm4add}. Beyond EER and minimum tandem detection cost
(min~t-DCF)~\cite{kinnunen2020tdcf}, we report bootstrap confidence intervals, expected
calibration error, Brier score, per-language and per-attack breakdowns, and
computational cost.

% [DUMMY] — headline numbers below are placeholders
The contributions of this work are as follows:
\begin{enumerate}
    \item \textbf{Architecture.} AuralGuard, a multi-view framework that couples an
        adaptive SSL foundation-model front-end (learnable layer attention with entropy
        regularization) with a vocoder-artifact branch through gated cross-attention
        fusion and a spectro-temporal graph back-end (Section~\ref{sec:method}).
    \item \textbf{Objective.} A one-class contrastive training objective (OC-Softmax
        with a supervised-contrastive auxiliary) that tightens the bona-fide manifold
        and improves rejection of unseen generators by up to 1.8 EER points over
        binary cross-entropy in ablation (Section~\ref{ssec:ablation}).
    \item \textbf{Robustness.} A channel-robustness augmentation curriculum that reduces
        EER degradation under codec compression by up to 40\% relative without
        sacrificing clean-condition accuracy (Section~\ref{ssec:robustness}).
    \item \textbf{Benchmark.} A seven-corpus, four-generation evaluation protocol
        covering zero-shot generalization, channel robustness, cross-lingual transfer,
        calibration, and explainability---including, to our knowledge, the first
        controlled comparison of specialized countermeasures against zero-shot audio
        large language models under identical cross-domain conditions
        (Section~\ref{sec:results}).
    \item \textbf{Open release.} Code, trained weights, evaluation manifests, an
        inference API, and an interactive demo are publicly released.
\end{enumerate}

The remainder of the paper is organized as follows. Section~\ref{sec:related} reviews
countermeasure architectures, speech foundation models, one-class learning, and the
emerging ALLM literature. Section~\ref{sec:method} details the AuralGuard architecture,
training objective, and augmentation curriculum. Section~\ref{sec:setup} describes
datasets, baselines, and implementation. Section~\ref{sec:results} presents in-domain,
zero-shot, robustness, cross-lingual, ablation, calibration, and efficiency results.
Section~\ref{sec:discussion} discusses limitations and ethical considerations, and
Section~\ref{sec:conclusion} concludes.
